\documentclass[11pt]{article}
% Кодировка и язык
\usepackage[utf8]{inputenc}           % Кодировка исходного файла (UTF-8)
\usepackage[T2A]{fontenc}             % Кодировка шрифтов для кириллицы
\usepackage[english,russian]{babel}   % Локализация: русские названия, переносы

% Поиск в PDF
\usepackage{cmap}

% Поля страницы
\usepackage{geometry}
\geometry{left=3.8cm, right=3.8cm, top=2.5cm, bottom=2.5cm, paperheight=29.7cm}

% Математика
\usepackage{amsmath, amssymb, amsthm} % Базовые пакеты для математики
\usepackage{mathtools}                % Расширение amsmath (доп. символы и окружения)

% Картинки
\usepackage{graphicx}                 % Вставка картинок (\includegraphics)
\graphicspath{{pictures/}}            % Путь к папке с картинками
\usepackage{float}                    % Позволяет использовать [H] для жёсткого позиционирования
\usepackage{wrapfig}                  % Обтекание текста вокруг картинок

% Таблицы
\usepackage{array}                    % Расширенные возможности для таблиц
\usepackage{booktabs}                 % Красивые таблицы (\toprule, \midrule, \bottomrule)
\usepackage{multirow}                 % Объединение строк в таблицах
\usepackage{tabularx}                 % Таблицы с автоматической шириной колонок
\usepackage{ragged2e}                 % Улучшенное выравнивание (используется в tabularx)

% Списки
\usepackage{enumitem}                 % Гибкая настройка списков (itemize, enumerate)

% Колонтитулы
\usepackage{fancyhdr}

% Гиперссылки
\usepackage{xcolor}                   % Работа с цветами
\usepackage{hyperref}                 % Гиперссылки, оглавление, перекрёстные ссылки

% Умные ссылки
\usepackage{cleveref}

% Типографика
\usepackage{microtype}                % Улучшает переносы и межсловные расстояния

% Подписи к картинкам и таблицам
\usepackage{caption}                  % Настройка подписей (шрифт, отступы)

% Рисование схем и графиков
\usepackage{tikz}                     % Рисование схем, графиков, диаграмм
\usepackage{pgfplots}                 % Построение графиков функций (на основе tikz)
\pgfplotsset{compat=newest}           % Версия совместимости pgfplots

% Алгоритмы и код
\usepackage{algorithm}                % Окружения для алгоритмов
\usepackage{algorithmic}              % Синтаксис для алгоритмов
\usepackage{listings}                 % Вставка кода с подсветкой синтаксиса

\usepackage{soul}                     % Выделение текста: подчёркивание, зачёркивание, подсветка
\usepackage{multicol}                 % Многоколоночная вёрстка текста
\usepackage{titlesec}                 % Настройка внешнего вида заголовков секций
\usepackage{parskip}                  % Отступы между абзацами вместо красной строки
\usepackage{setspace}                 % Настройка межстрочного интервала
\usepackage{tcolorbox}                % Цветные блоки для заметок, цитат и выделений

% Палитра Quiet Light из VS Code
\definecolor{quietLightForeground}{HTML}{333333}
\definecolor{quietLightEditorBg}{HTML}{F5F5F5}
\definecolor{quietLightPanelBg}{HTML}{F2F2F2}
\definecolor{quietLightBorder}{HTML}{C9D0D9}
\definecolor{quietLightComment}{HTML}{AAAAAA}
\definecolor{quietLightKeyword}{HTML}{4B69C6}
\definecolor{quietLightType}{HTML}{7A3E9D}
\definecolor{quietLightConstant}{HTML}{9C5D27}
\definecolor{quietLightString}{HTML}{448C27}
\definecolor{quietLightFunction}{HTML}{AA3731}
\definecolor{quietLightPunctuation}{HTML}{777777}

% Совместимость со старыми именами цветов
\colorlet{pblue}{quietLightKeyword}
\colorlet{pgreen}{quietLightString}
\colorlet{pred}{quietLightFunction}
\colorlet{pgrey}{quietLightPunctuation}

\onehalfspacing

\titleformat{\section}
  {\Large\bfseries}
  {\thesection}
  {0.7em}
  {}

\titleformat{\subsection}
  {\large\bfseries}
  {\thesubsection}
  {0.7em}
  {}

\lstset{
    basicstyle=\ttfamily\small\color{quietLightForeground},
    backgroundcolor=\color{quietLightEditorBg},
    frame=single,
    rulecolor=\color{quietLightBorder},
    breaklines=true,
    columns=fullflexible,
    keepspaces=true,
    showstringspaces=false,
    keywordstyle=\color{quietLightKeyword}\bfseries,
    morekeywords=[2]{define,elif,else,endif,error,if,ifdef,ifndef,include,line,pragma,undef,warning},
    keywordstyle=[2]\color{quietLightKeyword},
    deletekeywords={bool,char,double,float,int,long,short,signed,unsigned,void},
    morekeywords=[3]{bool,char,double,float,int,long,short,signed,unsigned,void,size_t,string},
    keywordstyle=[3]\color{quietLightType},
    identifierstyle=\color{quietLightType},
    commentstyle=\color{quietLightComment}\itshape,
    stringstyle=\color{quietLightString},
    numberstyle=\color{quietLightConstant},
}

\lstdefinestyle{nolint}{
    basicstyle=\ttfamily\small\color{quietLightForeground},
    backgroundcolor=\color{quietLightPanelBg},
    columns=fullflexible,
    keepspaces=true,
    breaklines=false,
    keywordstyle=\color{quietLightKeyword},
    commentstyle=\color{quietLightKeyword},
    stringstyle=\color{quietLightKeyword},
    numberstyle=\color{quietLightKeyword},
}

\newcommand{\nolint}{\lstinline[style=nolint]}

\newtcolorbox{habrquote}{
    colback=quietLightPanelBg,
    colframe=quietLightBorder,
    boxrule=0.4pt,
    arc=2pt,
    left=8pt,
    right=8pt,
    top=6pt,
    bottom=6pt
}

% Настройка нумерации секций и подсекций
\renewcommand{\thesection}{\arabic{section}}
\renewcommand{\thesubsection}{\thesection.\arabic{subsection}}


% Настройка списков: убираем лишние отступы слева
\setlist[itemize]{leftmargin=*}
\setlist[enumerate]{leftmargin=*}


% Определение окружений для теорем, лемм, определений и т.д.
% Нумерация привязана к номеру секции
\newtheorem{theorem}{Теорема}[section]
\newtheorem{definition}{Определение}[section]
\newtheorem{remark}{Замечание}[section]
\newtheorem{lemma}{Лемма}[section]
\newtheorem{example}{Пример}[section]
\newtheorem{consequence}{Следствие}[section]
\newtheorem{algorithmm}{Алгоритм}
\newtheorem{question}{Вопрос}[section]
\newtheorem{answer}{Ответ}[section]
\newtheorem{statement}{Утверждение}[section]


\DeclareMathOperator{\sign}{sign}
\DeclareMathOperator{\Imm}{Im}
\DeclareMathOperator{\rg}{rg}

% Настройка колонтитулов
\pagestyle{fancy}
\fancyhf{} % Очистить все поля
\fancyhead[L]{Кабы не было dynamic interpreter}
\fancyhead[R]{\href{https://t.me/ZNickA8871}{Автор}}
\fancyfoot[C]{\thepage}               % Справа внизу

\setlength{\headheight}{13.6pt}

\hypersetup{
    colorlinks=true,
    linkcolor=blue,
    citecolor=blue,
    filecolor=blue,
    urlcolor=blue,
}

\begin{document}

{\Huge\bfseries Лекция 2. \newlineДинамические библиотеки\par}

\vspace{2em}

\section{Зачем?}

Команда \nolint|ldd| для статически  слинкованных файлов (флаг -static при линковке) пишет ``not a dynamic executable''. Для нас это значит, что они запустятся незавсимо от того, есть ли у нашего пользователя те библиотеки, которые мы используем. Но при такой сборке даже простой hello world весит около 2 мегабайт. Это больно, особенно когда хочется иметь в системе много маленьких программ на все случаи жизни. Вот эту проблему и решают динамические библиотеки. Они уменьшают объём исполняемого файла,

Посмотрим, например, утилиту ls:
\newline
\nolint|which ls| покажет нам, где она находится,
\newline
\nolint|readelf -h <path>| сообщит, что она слинкована динамически,
\newline
\nolint|ldd <path>| выведет её зависимости.

Прикол в том, что эти зависимости загружаются в оперативную память один раз для всех, кому они нужныб и все запущенные программы просто переиспользуют один и тот же инстанс в памяти (точнее, переиспользуются только read-only секции, но об этом в следующих лекциях).

А теперь мы, допустим, мы захотели создать свою библиотеку.
Создадим два файла:

\nolint|multiply.cpp|, наша библиотека:
\begin{lstlisting}[language=C++]
int multiply(int x, int y) {
  return x * y;
}
\end{lstlisting}

\nolint|main.cpp|, который будет ее использовать:
\begin{lstlisting}[language=C++]
#include <iostream>

int multiply(int, int);

int main() {
  int x, y;
  std::cin >> x >> y;
  std::cout << multiply(x, y);
}
\end{lstlisting}

\newpage

Для начала соберём объектные файлы, а затем слинкуем их:

\begin{lstlisting}
g++ -c main.cpp
g++ -c multiply.cpp
g++ main.o multiply.o
\end{lstlisting}

И оно даже заработает. Но сейчас мы просто статически прилинковали multiply.o к main.o, а мы хотели именно динамическую библиотеку.

\subsection{Создание динамической библиотеки}

Соберём файл с расширением \nolint|.so|:

\begin{lstlisting}
  g++ -shared multiply.cpp -o multiply.so
\end{lstlisting}

Забавно, что у полученного файла может быть выставлен бит исполнения (\nolint|ls -l multiply.so|),
но попытка запустить его как обычную программу заканчивается ошибкой segmentation fault.
Посмотрим на файл подробнее:
\begin{lstlisting}
readelf -a multiply.so
\end{lstlisting}
В заголовке он обозначен как динамическая библиотека:
\begin{lstlisting}
Type: DYN (Shared object file)
\end{lstlisting}

В файле есть таблица динамических символов \nolint|.dynsym|. В ней находятся символы, которые может использовать динамический линковщик при запуске программы. Есть ещё похожая таблица \nolint|.symtab|, в ней содержатся все символы (и глобальные, и локальные), нужна она в основном для отладки и может быть удалена утилитой \nolint|strip|.

Утилита \nolint|objcopy| позволяет редактировать ELF-файлы. Ещё для работы с ними есть
полезная утилита \nolint|patchelf|.

Несмотря на то, что код динамической библиотеки загружается в оперативную память один раз для
нескольких процессов, секции \nolint|.data| и \nolint|.bss| у каждого процесса свои. Их страницы могут
физически быть общими до первой записи из-за оптимизации \textit{copy-on-write} (COW).

\subsection{Линковка с библиотекой}

Ключ \nolint|-l<name>| добавляет библиотеку при линковке:

\begin{lstlisting}
g++ main.cpp -lmultiply
\end{lstlisting}

Пока не работает: линковщик ожидает, что имя файла библиотеки начинается с префикса
\nolint|lib|, поэтому переименуем файл:

\begin{lstlisting}
mv multiply.so libmultiply.so
\end{lstlisting}

Всё ещё не работает: линковщик ищет библиотеки только в определённых каталогах, в том числе
в \nolint|/usr/lib|. Можно скопировать библиотеку туда:

\begin{lstlisting}
sudo cp libmultiply.so /usr/lib
g++ main.cpp -lmultiply
\end{lstlisting}

Теперь сборка проходит, а зависимости программы можно посмотреть с помощью \nolint|ldd|. Однако
копировать тестовую библиотеку в системный каталог необязательно. Удалим её оттуда и добавим
текущий каталог в путь поиска линковщика ключом \nolint|-L|:

\begin{lstlisting}
sudo rm /usr/lib/libmultiply.so
g++ main.cpp -lmultiply -L.
\end{lstlisting}

Сборка теперь пройдет, но программа при запуске может не найти библиотеку. Путь
поиска на этом этапе можно задать переменной окружения \nolint|LD_LIBRARY_PATH|, но на некоторых системах он уже задан как \nolint|/usr/local/lib:|, а пустой элемент после двоеточия означает текущую папку, поэтому, вероятно, программа запустится и так. Иначе, можно явно указать директорию поиска вот так:

\begin{lstlisting}
LD_LIBRARY_PATH=. ./a.out
\end{lstlisting}

Другой вариант --- записать путь поиска в сам исполняемый файл при линковке:

\begin{lstlisting}
g++ main.cpp -lmultiply -L. -Wl,-rpath=<absolute path>
g++ main.cpp -lmultiply -L. -Wl,-rpath=.
readelf -d a.out
\end{lstlisting}

В выводе \nolint|readelf| появится запись \nolint|RUNPATH|:

\begin{lstlisting}
0x000000000000001d (RUNPATH) Library runpath: [.]
\end{lstlisting}

Итак, ключ \nolint|-L| задаёт путь поиска библиотек во время сборки, а \nolint|LD_LIBRARY_PATH| и
\nolint|RUNPATH| --- во время запуска.

\subsection{SONAME}

В заголовке динамической библиотеки может быть указано поле \nolint|SONAME|.
Пусть при линковке используется файл \nolint|libmultiply.so|, а в его поле
\nolint|SONAME| записано другое имя, например, \nolint|libmultiplyhaha.so|. Тогда линковщик найдёт библиотеку по имени файла, но в список зависимостей
запишет именно \nolint|SONAME|.

Это позволяет при сборке использовать конкретную версию библиотеки, например
\nolint|libfoo.so.1.4.7|, а при запуске искать её по символической ссылке вроде
\nolint|libfoo.so.1|. Библиотеку можно обновлять, пока её интерфейс остаётся совместимым.
\newpage

Зададим нашей библиотеке другой \nolint|SONAME| и ещё раз соберём программу:

\begin{lstlisting}
g++ -shared multiply.cpp -o libmultiply.so -Wl,-soname=libmultiplyhaha.so
g++ main.cpp -lmultiply -L. -Wl,-rpath=.
ldd a.out
\end{lstlisting}

В результате в списке зависимостей окажется новое имя, но файла с таким именем ещё нет:

\begin{lstlisting}
linux-vdso.so.1
libmultiplyhaha.so => not found
libstdc++.so.6 => /lib/x86_64-linux-gnu/libstdc++.so.6
libc.so.6 => /lib/x86_64-linux-gnu/libc.so.6
libm.so.6 => /lib/x86_64-linux-gnu/libm.so.6
/lib64/ld-linux-x86-64.so.2
libgcc_s.so.1 => /lib/x86_64-linux-gnu/libgcc_s.so.1
\end{lstlisting}
После этого пересоберем библиотеку без \nolint|-soname|
\begin{lstlisting}
g++ -shared multiply.cpp -o libmultiply.so
\end{lstlisting}

\subsection{Patchelf}

Путь поиска можно добавить и в уже собранный ELF-файл. Сначала соберём программу
без \nolint|RUNPATH|, а затем изменим его утилитой \nolint|patchelf|:

\begin{lstlisting}
g++ main.cpp -lmultiply -L.
patchelf a.out --set-rpath .
\end{lstlisting}

После этого программа запускается. Врзможно, patchelf у вас не установлен, на Ubuntu ставим так:
\begin{lstlisting}
sudo apt install patchelf
\end{lstlisting}

\newpage
\section{Динамический линковщик}

Возьмем код из предыдущей лекции:
\begin{lstlisting}[language=C++]
#include <iostream>

int main() {
  int x;
  std::cin >> x;
  std::cout << x + 5;
  return 0;
}
\end{lstlisting}
Соберем объектный файл, а затем попробуем его слинковать. Чтобы было от чего отталкиваться, посмотрим зависимости обычного \nolint|a.out|:

\begin{lstlisting}
ldd a.out
\end{lstlisting}

\begin{lstlisting}[style=nolint]
linux-vdso.so.1
libstdc++.so.6 => /lib/x86_64-linux-gnu/libstdc++.so.6
libc.so.6 => /lib/x86_64-linux-gnu/libc.so.6
libm.so.6 => /lib/x86_64-linux-gnu/libm.so.6
/lib64/ld-linux-x86-64.so.2
libgcc_s.so.1 => /lib/x86_64-linux-gnu/libgcc_s.so.1
\end{lstlisting}

Попробуем вызвать \nolint|ld| вручную и перечислить ему нужные библиотеки:

\begin{lstlisting}
ld main.o /lib/x86_64-linux-gnu/libstdc++.so.6 \
  /lib/x86_64-linux-gnu/libc.so.6 \
  /lib/x86_64-linux-gnu/libm.so.6 \
  /lib64/ld-linux-x86-64.so.2 \
  /lib/x86_64-linux-gnu/libgcc_s.so.1
\end{lstlisting}
Получим warning
\begin{lstlisting}
ld: warning: cannot find entry symbol _start; defaulting to
\end{lstlisting}
А при запуске ошибку
\begin{lstlisting}
bash: ./a.out: cannot execute: required file not found
\end{lstlisting}
\newpage
Сначала разберёмся с warningом. Добавим объектные файлы, с которых начинается и которыми заканчивается наша программа:

\begin{lstlisting}
ld main.o \
  /lib/x86_64-linux-gnu/libstdc++.so.6 \
  /lib/x86_64-linux-gnu/libm.so.6 \
  /lib64/ld-linux-x86-64.so.2 \
  /lib/x86_64-linux-gnu/libgcc_s.so.1 \
  /lib/x86_64-linux-gnu/libc.so.6 \
  /usr/lib/gcc/x86_64-linux-gnu/14/crtendS.o \
  /usr/lib/x86_64-linux-gnu/crtn.o \
  /usr/lib/x86_64-linux-gnu/Scrt1.o \
  /usr/lib/x86_64-linux-gnu/crti.o \
  /usr/lib/gcc/x86_64-linux-gnu/14/crtbeginS.o
\end{lstlisting}
Теперь она даже соберется, но тут всё ещё есть две ошибки.
При ручной линковке можно получить другой путь к динамическому линковщику, из-за чего
при запуске возникнет ошибка \nolint|No such file or directory| или упомянутая ранее \nolint|bash: ./a.out: cannot execute: required file not found|. Путь к интерпретатору можно посмотреть
двумя способами:

\begin{lstlisting}
patchelf --print-interpreter <file>
readelf -l <file>
\end{lstlisting}

А явно задать его линковщику можно ключом:

\begin{lstlisting}
-dynamic-linker /lib64/ld-linux-x86-64.so.2
\end{lstlisting}

Попробуем:
\begin{lstlisting}
ld main.o -dynamic-linker /lib64/ld-linux-x86-64.so.2 \
  /lib/x86_64-linux-gnu/libstdc++.so.6 \
  /lib/x86_64-linux-gnu/libm.so.6 \
  /lib64/ld-linux-x86-64.so.2 \
  /lib/x86_64-linux-gnu/libgcc_s.so.1 \
  /lib/x86_64-linux-gnu/libc.so.6 \
  /usr/lib/gcc/x86_64-linux-gnu/14/crtendS.o \
  /usr/lib/x86_64-linux-gnu/crtn.o \
  /usr/lib/x86_64-linux-gnu/Scrt1.o \
  /usr/lib/x86_64-linux-gnu/crti.o \
  /usr/lib/gcc/x86_64-linux-gnu/14/crtbeginS.o
\end{lstlisting}
...
\newpage
И с высокой вероятностью получим segfault.
Вторая наша ошибка была в том, что мы указали файлы не в том порядке. Попробуем теперь так:
\begin{lstlisting}
ld -dynamic-linker /lib64/ld-linux-x86-64.so.2 \
  /usr/lib/x86_64-linux-gnu/Scrt1.o \
  /usr/lib/x86_64-linux-gnu/crti.o \
  /usr/lib/gcc/x86_64-linux-gnu/14/crtbeginS.o \
  main.o \
  /lib/x86_64-linux-gnu/libstdc++.so.6 \
  /lib/x86_64-linux-gnu/libm.so.6 \
  /lib64/ld-linux-x86-64.so.2 \
  /lib/x86_64-linux-gnu/libgcc_s.so.1 \
  /lib/x86_64-linux-gnu/libc.so.6 \
  /usr/lib/gcc/x86_64-linux-gnu/14/crtendS.o \
  /usr/lib/x86_64-linux-gnu/crtn.o
\end{lstlisting}
Иии... Всё заработает.

Отдельный интересный вопрос --- как запускается сам динамический линковщик, который тоже
является \nolint|.so|-файлом, но при этом должен запустить нашу программу. Прикол в том, что у
\nolint|/lib64/ld-linux-x86-64.so.2|, и у \nolint|/lib/x86_64-linux-gnu/libc.so.6| есть точка входа: их можно использовать как
библиотеки, а можно запускать напрямую.

\section{Секции ELF и память}

Подробную информацию о секциях можно получить командой:

\begin{lstlisting}
readelf -t a.out
\end{lstlisting}

Основные секции:

\begin{itemize}
    \item \nolint|.text| --- машинный код;
    \item \nolint|.data| --- изменяемые инициализированные глобальные и статические данные;
    \item \nolint|.rodata| --- неизменяемые инициализированные глобальные и статические данные;
    \item \nolint|.bss| --- неинициализированные или инициализированные нулём глобальные и статические
    данные (\textit{Block Started by Symbol}).
\end{itemize}

Если привести указатель на изначально константный объект к неконстантному типу и
попытаться выполнить запись, поведение программы будет неопределённым. Для глобального или
статического объекта, попавшего в доступную только для чтения память, на практике это обычно приведёт
к segmentation fault. Константная локальная переменная на стеке обычно не вызовет такой ошибки, но
запись в неё всё равно остаётся неопределённым поведением.

Команда \nolint|size <file>| выводит размеры основных секций бинарного файла.

P.S. Часть примеров здесь следует собирать как PIE, но разбор этой темы слишком объемный для нашей лекции, а примеры будут работать и так, т.к. они довольно простые. Желающим разобраться предлагается воспользоваться гуглом и гпт.
\end{document}
